% =============================================================================
% Kapitel 3: Methodik
% =============================================================================

\section{Methodik}
\label{sec:methodik}

\subsection{Untersuchungsansatz}
\label{sec:methodik-ansatz}

Die Untersuchung folgt dem Prinzip der \textit{Fault Injection}: dem gezielten Einbringen
definierter Fehlerzustände in ein laufendes System, um dessen Reaktion empirisch zu
beobachten \cite{hsueh1997}. Für jeden Failure Case wird eine Hypothese über das erwartete
Systemverhalten formuliert, ein reproduzierbares Fehlerszenario hergestellt und das
tatsächliche Verhalten anhand definierter Messgrößen erfasst. Wo ein Fehlerzustand ohne
Eingriff in den Quellcode herstellbar ist, werden ausschließlich POSIX-Signale oder
externe Werkzeuge eingesetzt; wo ein Eingriff nötig ist, wird dieser minimal gehalten und
nach der Messung automatisch zurückgenommen.

\subsection{Testumgebung}
\label{sec:methodik-umgebung}

Alle Messungen erfolgen am Referenzsystem im Labor Bad Hersfeld. Das System wird über
einen Simulator betrieben, der die Kamerahardware durch das Einspielen vorab erfasster
Bilddaten ersetzt. Sämtliche Failure Cases sind dadurch ohne echte Kamerahardware unter
kontrollierten Bedingungen reproduzierbar. Tabelle~\ref{tab:testumgebung} fasst die
relevanten Eigenschaften zusammen.

\begin{table}[htbp]
    \centering
    \caption{Eigenschaften der Testumgebung}
    \label{tab:testumgebung}
    \begin{tabularx}{\textwidth}{lX}
        \hline
        \textbf{Eigenschaft} & \textbf{Wert} \\
        \hline
        Betriebssystem      & Linux 6.17 (x86\_64) \\
        Arbeitsspeicher     & 64 GB \\
        GPU                 & NVIDIA GeForce RTX 5090 (32 GB VRAM) \\
        Laufzeitumgebung    & Python 3.13 \\
        Startskript         & \texttt{main.py} (mit geänderter Konfiguration für Simulation) \\
        Bilddatenquelle     & 462 Bilder = 154 Inspektionen \\
        \hline
    \end{tabularx}
\end{table}

\subsection{Referenzmessung (Baseline)}
\label{sec:methodik-baseline}

Vor der Fehlerinjektion wurde eine Referenzmessung im Normalbetrieb ohne injizierte Fehler
durchgeführt. Sie dient als Vergleichsbasis für alle Failure Cases. Das System startet
neben dem Hauptprozess vier Service-Prozesse für \texttt{Logger}, \texttt{FrameGrabber},
\texttt{Inspection\allowbreak{}Data\allowbreak{}Handler} und \texttt{PipelineManager} sowie
einen Pool von 12 Worker-Prozessen. Im stationären Zustand erreicht es einen Durchsatz von
1,74~Inspektionen pro Sekunde bei einer mittleren Pipeline-Laufzeit von 383~ms. Die
zentrale Log-Queue bleibt dabei nahezu leer.

\subsection{Messgrößen}
\label{sec:methodik-messgroessen}

Für jeden Case werden vergleichbare Messgrößen erhoben:

\begin{itemize}
    \item \textit{Time-to-Detection (TTD):} Zeit von der Fehlerinjektion bis zur
          Erkennung durch das System (Shutdown oder Statusmeldung).
    \item \textit{Folgeverhalten:} qualitative Beschreibung der Reaktion der übrigen
          Prozesse (Weiterlauf, Deadlock, Leerschleife, Absturz).
    \item \textit{Queue-Tiefe und Speicherfüllstand über die Zeit:} quantitative
          Erfassung schleichender und kaskadierender Effekte.
    \item \textit{Graceful Degradation:} ob das System bei einem Teilausfall in einen
          kontrollierten Zustand übergeht oder vollständig ausfällt.
\end{itemize}
